\NewDocumentCommand{\getelem}{O{\data} m m m}{%m - mandatory -- обов'язковий параметр
	\pgfplotstablegetelem{#2}{#3}\of#1
	\pgfmathsetmacro{#4}{\pgfplotsretval}
}

%------------------Функція для оформлення змінної та її похибки похибок
%перед записуванням даних слід додати \pgfkeys{/pgf/fpu, /pgf/fpu/output format=fixed}, особливо якщо похибка 1000000%
% після слід додати \pgfkeys{/pgf/fpu=false}
\NewDocumentCommand{\printwitherror}{m m}{%#1 a #2 aerr 
	\pgfmathfloatparsenumber{#1}
	\pgfmathfloattomacro{\pgfmathresult}{\FVal}{\ManVal}{\ExVal}
	\pgfmathfloatparsenumber{#2}
	\pgfmathfloattomacro{\pgfmathresult}{\FErr}{\ManErr}{\ExErr}
	\pgfmathtruncatemacro{\pr}{ifthenelse(int(\ExVal-\ExErr+1)>0,int(\ExVal-\ExErr+1),int(0)}%Знаходимо точність відображення значення беручи до уваги лише один знак після крапки в похибці
	%\pgfmathtruncatemacro замість \pgfmathsetmacro щоб відкинути зайві нулі бо інакше вони вивидуться в результат. Наприклад, \pgfkeys{/pgf/number format/.cd, precision=4.00000, zerofill, fixed} \pgfmathprintnumber{1.23456} видасть .000001.2346
	%Використання int тут дуже важливе інакше макрос не коректно передаєтся в percision
	\pgfkeys{/pgf/number format/.cd, precision=\pr, fixed zerofill,fixed}
	(\pgfmathprintnumber{\ManVal}\pm
	\pgfmathparse{\ManErr/10^(\ExVal-\ExErr)}
	\pgfmathprintnumber{\pgfmathresult})
	\ifnum \ExVal=0
	\relax 
	\else
	\cdot 10^{\pgfmathprintnumber[zerofill=false,precision=0]{\ExVal}}
	\fi
}
\NewDocumentCommand{\printinDMS}{m}{
	\pgfkeys{/pgf/fpu, /pgf/number format/.cd, precision=0, fixed}
	\pgfmathfloatparsenumber{#1}
	\pgfmathsetmacro{\DegresNR}{\pgfmathresult}
	\pgfmathsetmacro{\Degres}{int(\DegresNR)}
	\pgfmathsetmacro{\MinutsNR}{(\DegresNR-\Degres)*60}
	\pgfmathsetmacro{\Minuts}{int(\MinutsNR)}
	\pgfmathsetmacro{\Seconds}{int((\MinutsNR-\Minuts)*60)}
	\pgfmathprintnumber{\Degres}^{\circ} \pgfmathprintnumber{\Minuts}^ {\prime}\pgfmathprintnumber{\Seconds}^{\prime\prime}
	\pgfkeys{/pgf/fpu=false}
}
